\documentclass[11pt]{article}

\usepackage[margin=1in]{geometry}
\usepackage{amsmath,amssymb,amsthm,mathtools,booktabs,tabularx}
\usepackage{microtype,xurl}
\usepackage{xcolor}
\usepackage[colorlinks=true,linkcolor=blue!55!black,
            citecolor=blue!55!black,urlcolor=blue!55!black]{hyperref}

\newtheorem{theorem}{Theorem}[section]
\newtheorem{proposition}[theorem]{Proposition}
\newtheorem{corollary}[theorem]{Corollary}
\theoremstyle{definition}
\newtheorem{definition}[theorem]{Definition}
\theoremstyle{remark}
\newtheorem{remark}[theorem]{Remark}

\newcommand{\dd}{\mathrm d}
\newcommand{\cert}[1]{\nolinkurl{#1}}
\newcolumntype{Y}{>{\raggedright\arraybackslash}X}

\title{The Gauge Algebra of Conformal Gravity\\
\large A Direct Check on the Bach Noether Identities}
\author{Claude Opus 5\thanks{Anthropic model; computations and drafting.}
\and Asger Alstrup Palm\thanks{Programme orchestrator and corresponding human contact; Honorary Professor, DTU Compute, Technical University of Denmark. \texttt{asger@area9.dk}.}}
\date{Working draft, 7 August 2026}

\begin{document}
\maketitle

\begin{abstract}
The Weyl action \(S=\alpha\int\!\sqrt{-g}\,C^{2}\) is invariant under
diffeomorphisms and under local Weyl rescalings.  That much is a calculation.
Whether \(\mathrm{Diff}\ltimes\mathrm{Weyl}\) is \emph{all} of its gauge
algebra is a different question, and a harder one, since it asks about
symmetries the action was not built to have.

This question has an answer, and it is not ours.  Boulanger and Henneaux
\cite{BoulangerHenneaux2001} showed by cohomological deformation theory that
the Weyl action is the only consistent deformation of linearized conformal
gravity with no more than four derivatives of the metric, and---the part that
bears on us---that diffeomorphism invariance \emph{follows from consistency}
rather than being imposed.  On the question of principle we add nothing.

What we add is a direct verification by an unrelated route.  By Noether's
second theorem gauge generators and Noether identities are the same data, so
the question becomes whether the Bach tensor satisfies any local identity
beyond \(\nabla^{a}B_{ab}=0\) and \(g^{ab}B_{ab}=0\)---a question about the
\emph{full nonlinear} action, asked order by order in the gauge parameter,
rather than about deformations of the free theory.  Working over a metric family carrying symbolic
parameters, so that a vanishing result is zero \emph{as a polynomial} rather
than at sampled fixtures, we find no further generator at derivative order
zero, one or three in the gauge parameter, at the curvature degrees stated.
Orders two and four are empty for a reason that needs no computation at all:
index parity.  Order three is the delicate case, because the third covariant
derivative of the Bach tensor is seventh order in the metric and most of an
affordable jet is truncation artefact; we measure how much is usable rather
than assume it, and settle the case by exhibiting a rank that does not move as
the usable range triples.

Two features of the problem make a naive count wrong, and both bite.  Gauge
symmetries form a module over \emph{functions}, so a generator multiplied by a
scalar is not a second generator; we find the trap is worse than that, since
\(g_{\mu\nu}(\nabla\!\cdot\!\xi)\) is the Weyl generator under a
\emph{differential} reparametrisation of its parameter, so the quotient must
be by differential operators and not merely by multiplication.  And every
action has trivial symmetries proportional to its own equations of motion; we
exclude these by construction rather than by assertion.

The result is a lower bound on completeness at the order computed, and is
reported as one.  We do not claim the algebra is closed at all orders.  We
claim that nothing hides below the order we reached, and we say exactly where
that order stops.
\end{abstract}

\section{The question}
\label{sec:question}

A gauge symmetry of an action is not a property one may simply declare.  The
Weyl action was constructed to be invariant under
\(g_{\mu\nu}\mapsto\Omega^{2}g_{\mu\nu}\) and under coordinate changes, and
exhibiting those invariances is routine.  What is not routine is the converse:
\emph{is there anything else?}

The distinction matters for the same reason it matters anywhere in a
reverse-physics ledger.  If \(\mathrm{Diff}\ltimes\mathrm{Weyl}\) is the whole
algebra, then it is a \emph{consequence} of the action and not an independent
entry in the list of things conformal gravity assumes.  If it is not, then
conformal gravity has a local symmetry nobody has written down, which would be
a considerably larger claim than the expected answer and would need
correspondingly better evidence.

\paragraph{This has been answered, and the answer is not ours.}
Boulanger and Henneaux \cite{BoulangerHenneaux2001} classified the consistent
deformations of \emph{linearized} conformal gravity in \(D=4\) and found that
with no more than four derivatives of the metric the only possibility is the
Weyl action---and that diffeomorphism invariance is not an input but a
consequence: ``No a priori requirement of invariance under diffeomorphisms is
imposed: this follows automatically from consistency.''  On the question of
principle raised above, that is the answer, and it predates this work by a
quarter century.

Two things nevertheless remain to be done, and they are what this paper does.
Their argument bounds the \emph{action} by its derivative order and proceeds
by deforming the free theory; ours bounds the \emph{gauge parameter} by its
derivative order and interrogates the full nonlinear action directly, through
its Noether identities.  The two routes share no machinery, so agreement
between them is a genuine cross-check rather than a restatement.  And the
order-by-order form makes explicit \emph{where} the ceiling sits, which a
deformation argument does not report in those terms.

The contribution is therefore a confirmation and a bound, not a discovery.  A
reader who wants to know whether \(\mathrm{Diff}\ltimes\mathrm{Weyl}\) is
forced should cite \cite{BoulangerHenneaux2001}; a reader who wants it checked
against the nonlinear action by exact computation may cite this.

\paragraph{Conventions.}
\(D=4\), signature \((-,+,+,+)\).  \(B_{\mu\nu}\) is the Bach tensor, and
\(E^{\mu\nu}\) the Euler--Lagrange expression of \(S\), so that
\(\delta S=\int\!E^{\mu\nu}\delta g_{\mu\nu}\); for the Weyl action these
agree up to a constant.  All arithmetic is exact over \(\mathbb{Q}\).

\section{Why it is computable}
\label{sec:computable}

Noether's second theorem identifies gauge generators with Noether identities.
A local operator \(\Delta_{\mu\nu}[\varepsilon]\), linear in a parameter
\(\varepsilon\), generates a gauge symmetry exactly when
\(E^{\mu\nu}\Delta_{\mu\nu}[\varepsilon]\) is a total divergence for every
metric --- equivalently when the \emph{adjoint} \(\Delta^{\dagger}[E]\)
vanishes identically.\footnote{The adjoint \emph{reverses} the order of the
derivatives: the outermost derivative acting on the parameter in
\(\Delta\) becomes the innermost acting on \(E\).  Since the derivatives do
not commute, assigning a candidate to a generator before computing the adjoint
is a reliable way to turn an ordering slip into a result.  We compute the
contractions first and interpret afterwards.}  So the question of
\S\ref{sec:question} becomes:

\begin{quote}
does \(B_{\mu\nu}\) satisfy any local identity beyond
\(\nabla^{\mu}B_{\mu\nu}=0\) and \(g^{\mu\nu}B_{\mu\nu}=0\)?
\end{quote}

The two known identities are the Noether content of the two known generators:
the divergence is the diffeomorphism identity, the trace is the Weyl identity.
The question is whether the list is complete.

The problem splits by the \emph{type} of the gauge parameter.  A scalar
parameter \(\sigma\) gives candidates \(\Delta_{\mu\nu}=T_{\mu\nu}\sigma\); a
vector parameter \(\xi_{\lambda}\) gives candidates carrying derivatives,
\[
  \Delta_{\mu\nu}[\xi] \;=\; S_{\mu\nu}{}^{\rho\sigma\tau\lambda}
  \nabla_{\rho}\nabla_{\sigma}\nabla_{\tau}\xi_{\lambda}
\]
and its lower-order analogues.  These are computed separately and the results
combined in \S\ref{sec:result}.

\paragraph{What ``over a family'' means here.}
Every identity below is evaluated on a metric family carrying symbolic
parameters, so a result of zero is zero \emph{as a polynomial in those
parameters} --- true for every member at once, not at sampled fixtures.  The
family is general in the metric (ten parameters span every component of a
symmetric \(4\times4\)) but each component carries a fixed polynomial pattern
in the coordinates.  It is a generic family near flat space, not every metric,
and \S\ref{sec:not} says so again.

\section{Two ways to overcount, both of which bite}
\label{sec:traps}

\subsection{The module is over functions, and the trap is worse than that}

Gauge symmetries form a module over functions of the fields, not a vector
space.  If \(\Delta_{\mu\nu}=g_{\mu\nu}\sigma\) is the Weyl generator, then so
is \(\Delta_{\mu\nu}=g_{\mu\nu}R\sigma\), and it is not a second symmetry ---
it is the first with its parameter reparametrised, \(\sigma\mapsto R\sigma\).
A kernel dimension that counts these separately overcounts, and any reported
count that does not say how it handled this should not be believed.

The vector sector shows the trap is strictly worse.  The degree-zero candidate
\(S_{\mu\nu}{}^{\rho\sigma}=g_{\mu\nu}g^{\rho\sigma}\) gives
\[
  \Delta_{\mu\nu} \;=\; g_{\mu\nu}\,(\nabla\!\cdot\!\xi),
\]
which is the Weyl generator with \(\sigma=\nabla\!\cdot\!\xi\): a
\emph{differential} reparametrisation of the parameter, not a functional one.
So the quotient must be taken by differential operators acting on the
parameter, and a count that quotients only by multiplication by functions
\emph{still} overcounts.  Concretely, the naive degree-zero kernel in the
vector sector has three elements and the generator count is two.

\subsection{Trivial symmetries, excluded by construction}

Every action, without exception, admits symmetries of the form
\(\Delta_{\mu\nu}=M_{\mu\nu\alpha\beta}E^{\alpha\beta}\) with \(M\)
antisymmetric under exchange of its index pairs.  They are infinite in number
and carry no information; a count that includes them is meaningless.

They are excluded here by construction rather than by assertion.  Every
coefficient enumerated is built from the metric and the curvature and never
from \(E\), and the Bach tensor is fourth order in the metric, so no candidate
at the degrees considered can contain it.

\section{The result}
\label{sec:result}

\begin{theorem}[Scalar sector]
\label{thm:scalar}
Among candidates \(\Delta_{\mu\nu}=T_{\mu\nu}\sigma\) with \(T\) built from the
metric and curvature at curvature degree at most two, there is exactly
\textbf{one} generator, the Weyl generator.  The kernel consists precisely of
the multiples \(f\,g_{\mu\nu}\), which are reparametrisations of it; every
candidate that is not such a multiple has nonvanishing adjoint.
\end{theorem}

The candidate list for Theorem~\ref{thm:scalar} has ten elements, not the seven
of the list we inherited.  Degree two has seven members in full --- three
scalars times \(g\), one scalar times \(\mathrm{Ric}\), and three carrying both
free indices on curvature --- and with degrees zero and one that is ten.  The
three that a \(g\)-and-\(\mathrm{Ric}\) list omits are exactly the ones
contracting against the full Riemann tensor, which is the shape that could
carry a second generator, since \(E^{\mu\nu}\) here is the Bach tensor and sees
the whole Weyl tensor rather than only its traces.

\begin{theorem}[Vector sector]
\label{thm:vector}
Among candidates with a vector parameter at derivative order \(k\le 3\), the
only generator is the diffeomorphism generator.  Specifically:
\begin{itemize}\itemsep2pt
\item \(k=0\) and \(k=1\): no further generator, at curvature degree up to two
  in the coefficient.
\item \(k=2\) and \(k=4\): \textbf{empty by index parity}, with no computation
  required.  A coefficient at order \(k\) carries \(2+k+1\) indices, and every
  tensor built from \(g\), \(g^{-1}\), \(\delta\) and undifferentiated
  curvature has even rank, so the odd-rank coefficient the even orders need
  cannot be built.
\item \(k=3\), curvature degree zero: no further generator, established as
  described in Remark~\ref{rem:truncation}---a rank stable under increasing
  resolution, with the rank deficit accounted for by a coincidence rather than
  a relation.  At curvature degree \emph{one} the analogous sweep is
  \textbf{open}: its kernel numbers are not yet interpretable, because the
  quotient by trivial relations is currently only by pairwise equality and a
  scalar multiple is as much an artefact as a duplicate.
\end{itemize}
\end{theorem}

\begin{remark}[The parity argument has a limit, which is why \(k=0\) was computed]
Derivatives of curvature have \emph{odd} rank --- \(\nabla_{a}R\) is rank one,
\(\nabla_{a}R_{bc}\) rank three --- so order zero is reachable after all, with
one derivative on the curvature in the coefficient.  It was therefore
enumerated rather than dismissed by the parity argument that disposes of orders
two and four.  Order zero is also the cheap sector: its adjoint is a
contraction rather than a divergence, with no integration by parts.
\end{remark}

\begin{remark}[Scope of the degree-one sweep]
\label{rem:shape}
Every candidate in that sweep has the form \(Y=F_{1}[\text{pair}]\times
F_{2}[\text{pair}]\) with \emph{both factors of rank two}, and the enumeration
is over which two of the four slots each factor occupies.  It is therefore
complete in \emph{placement} within that shape class, and \emph{not} complete
in shape: coefficients that do not factor into two rank-two pieces---for
instance Riemann carrying three free indices, \(R_{\rho\sigma\tau
a}B^{a}{}_{\lambda}\)---are outside it.  At curvature degree one the result is
consequently weaker than at degree zero, and is stated as such.
\end{remark}

\begin{remark}[How order three is settled, and what nearly went wrong]
\label{rem:truncation}
\(\nabla^{3}B\) costs seven derivatives, so most of an affordable jet is
truncation artefact.  Rather than trust a budget we \emph{measure} the usable
range: the contractions whose innermost operation is \(\nabla\!\cdot\!B\)
must vanish identically, and the highest order at which they still do says how
far the data is good.  The arithmetic predicted no usable orders; the
measurement found three at jet degree seven, and the range grows with the jet
degree---orders \(0\ldots2,3,4,6\) at degrees \(7,8,9,10\).

Two facts then settle the case.  The rank of the surviving contractions is
\(2\) of \(3\) \emph{at every one of those degrees}: a number that does not
move while the usable data triples is not a resolution artefact.  And the
deficit is a \emph{coincidence}---exactly one pair among the three yields the
same adjoint up to sign, three vectors with one coinciding pair have rank at
most two, and the rank is two.  The coincidence accounts for the whole
deficit, leaving nothing to be a symmetry.

An earlier version reached this conclusion by a route that does not survive
scrutiny: it exhibited a candidate relation and found it failed when compared
over the \emph{whole} jet, but the failure lay entirely above the valid range.
Right answer, invalid reason.  We record it because the failure mode is
general---``testing on the full jets'' means nothing where the jet carries no
information---and because the repair, measuring the usable range against a
known answer before comparing anything, is the transferable part.
\end{remark}

\begin{corollary}
At derivative orders zero and one (curvature degree \(\le2\)) and at order
three (curvature degree zero), the gauge algebra of
\(S=\alpha\int\!\sqrt{-g}\,C^{2}\) is
\(\mathrm{Diff}\ltimes\mathrm{Weyl}\); orders two and four are empty by
parity.  This agrees with what \cite{BoulangerHenneaux2001} establishes by
deformation theory, reached independently, over the range where our
computation resolves.
\end{corollary}

\section{The negatives are the result}
\label{sec:negatives}

The positive cases in Theorems~\ref{thm:scalar} and \ref{thm:vector} are
already known and are recomputed only as premises: the diffeomorphism adjoint
is \(\nabla^{\mu}E_{\mu\nu}\) and the Weyl adjoint is the trace.  What makes
``exactly these'' evidence rather than an enumeration that never ran is that
curvature-weighted candidates \emph{must fail}, and do.

At order one, three such candidates fail: \(R\) weighting the diffeomorphism
candidate, \(\mathrm{Ric}\) contracted into the trace slot --- which requires
\(B^{\mu\nu}R_{\mu\nu}\equiv0\), and that is not an identity --- and
\(\mathrm{Ric}\) as a mixed factor.  At order three and curvature degree one,
all eight sub-families fail.

Two structural controls carry more weight than the negatives, because they are
known answers rather than expectations.  First, whenever the innermost
derivative of a third-order candidate contracts with an index of \(B\), the
innermost operation is \(\nabla\!\cdot\!B\), which is zero; those contractions
must vanish identically, and do, through a rank-five object.  That is the
diffeomorphism identity differentiated twice, and if it failed nothing else
would mean anything.  Second, the contraction of Riemann's second index pair
with a symmetric tensor is zero \emph{before any derivative is taken}, because
Riemann is antisymmetric there and \(B\) is symmetric; that tests the
construction rather than the adjoint machinery.

\begin{remark}[One control is vacuous and is reported as such]
The degree-zero Weyl candidate has adjoint \(\nabla^{\sigma}(\operatorname{tr}
E)\).  The Bach tensor is traceless, so this vanishes \emph{for that reason}
and not independently.  It reduces to the known trace identity; it is not a
check, and a control evaluated where the thing it measures degenerates proves
nothing.
\end{remark}

\subsection*{A complementary negative, arrived at independently}

While this work was in preparation, Hobson, Barker and
Lasenby~\cite{HobsonBarkerLasenby2026} extended the Bessel--Hagen
construction to the conformal group and proved that linearised Weyl-squared
gravity admits \emph{no} strict local, polynomial, symmetric,
dimension-four rank-two tensor quadratic in \(h_{\mu\nu}\) that is both gauge
invariant and conserved on the Bach shell---so no strictly gauge-invariant
local energy--momentum tensor exists there.

That is a different object in a different regime: a conserved \emph{current}
in the linearised theory, where ours is a gauge \emph{generator} in the full
nonlinear one.  Neither result implies the other.  But they are the same kind
of statement about the same theory---an exhaustive search that comes back
empty---reached by unrelated methods, and we record the pairing because it is
the more informative reading of both.  Weyl-squared gravity appears to be a
theory whose structure is pinned as much by what it does not admit as by what
it does.

\section{Rank is a floor, and the data has an edge}
\label{sec:rank}

The adjoint of a candidate is a linear combination of finitely many
contractions, so the symmetries are that combination's kernel.  Two readings of
such a computation are wrong and both were live here.

Showing each contraction separately nonvanishing does \emph{not} show they are
independent: a combination of nonvanishing objects can vanish, and that
combination would be a symmetry.  Conversely, a rank read at a base point is
only a \emph{lower} bound on the generic rank, so a rank deficit there is not a
relation.  What settles it is exhibiting the candidate combination --- by
carrying an identity through the elimination, so a row whose left half reduces
to zero returns its own coefficients --- and then testing that combination on
the full jets, where a genuine relation must be the zero polynomial at every
order.

This procedure has an edge, and at order three we ran off it.  ``Testing on
the full jets'' is only meaningful where the jet carries information, and after
seven derivatives most of a degree-seven jet does not.  A witness that
\emph{fails} above the valid range has not failed at all.  That is how the
withdrawn conclusion of Remark~\ref{rem:truncation} arose, and the guard
against it is to measure the usable range against a known answer before
comparing anything---which is now done, and reported alongside the rank.

\begin{remark}[A width, not just an order]
A relation may verify identically over a narrow family and still be false of
the theory.  One did: over a one-parameter family a relation among the
\(R\)-weighted order-three candidates verified at every order, and widening the
family to two parameters removed it.  It was an artefact of a family too thin
to separate the candidates.  The parameter count is a knob to vary, not to fix
at the cheapest value that runs, and the computations reported here are at the
wider setting.
\end{remark}

\section{What this paper does not establish}
\label{sec:not}

\begin{itemize}\itemsep2pt
\item \textbf{Not all orders, and not order three at curvature degree one.}
  The result covers derivative orders zero and one, order three at curvature
  degree zero, and two and four by parity.  Order three at curvature degree
  one is open---its kernel numbers are not yet interpretable, since the
  quotient by trivial relations is only by pairwise equality.  Order five and
  above is untouched.
\item \textbf{Not all metrics.}  The family is general in the metric but each
  component carries a fixed polynomial pattern in the coordinates.  This is a
  generic family near flat space.  An identity that fails only on metrics far
  from this family would not be seen.
\item \textbf{Not a proof that the ranks are full.}  Only individual witnesses
  were tested against the full jets.  A \emph{combination} of witnesses could
  vanish while none does alone, so ``no witness verified'' is strong evidence
  for full rank rather than a proof of it.
\item \textbf{Not complete in shape at curvature degree one.}  See
  Remark~\ref{rem:shape}: that sweep is exhaustive in placement within a
  restricted shape class, not over all coefficients of that degree.
\item \textbf{No principled stopping rule for the family width.}  The family
  carries two parameters because one produced a spurious relation
  (\S\ref{sec:rank}) and two did not.  That is a repair, not a criterion:
  nothing here shows two is \emph{enough}, and we do not have an argument that
  the family spans the jet directions the identities probe.  It does not.
\item \textbf{One implementation.}  Unlike some results in this programme,
  nothing here carries an independent second rail.  A systematic error in the
  shared curvature layer would not be caught by these checks, though the
  Riemann comparison in \S\ref{sec:verification} narrows that exposure.
\item \textbf{\(D=4\) and one signature.}  Nothing is claimed in other
  dimensions, where the counting arguments---index parity in particular---would
  have to be redone.
\item \textbf{No novelty in the symmetries themselves.}  That the Weyl action
  is diffeomorphism and Weyl invariant is classical.  The contribution is the
  converse direction --- that at bounded order there is nothing else --- which
  is what converts an assumption into a consequence.
\item \textbf{Nothing about the quantum theory.}  Everything here is a
  classical local computation.  Whether these symmetries survive quantisation
  is a separate question with its own obstructions, treated elsewhere in this
  programme.
\end{itemize}

\section{Verification}
\label{sec:verification}

Each result is a fail-closed certificate with a deterministic verifier, computed
in Forge over exact rationals with nested jets: the outer degree carries the
coordinate dependence and the inner degree the family parameters.  Both
truncation degrees are guarded, since a saturated jet compares truncation caps
rather than objects and an exact-identity test between saturated jets is
inconclusive rather than false.

The euler gate stands at 9 of 10 by design: its tenth check extends the same
machinery to the conformally invariant subspace, does not pass its own
known-answer control, and is therefore withheld.  A gate that reported 10 of 10
while that control failed would be the more worrying number.

\begin{center}
\begin{tabularx}{\textwidth}{@{}lYl@{}}
\toprule
result & certificate & rail \\
\midrule
Thm.~\ref{thm:scalar} & \cert{REVERSE\_PHYSICS\_NOETHER\_IDENTITIES\_V1}   & euler gate, 9/10 \\
Thm.~\ref{thm:vector} & \cert{REVERSE\_PHYSICS\_GAUGE\_VECTOR\_SECTOR\_V1} & family gate, 18/18 \\
curvature inputs      & \cert{REVERSE\_PHYSICS\_SYMBOLIC\_FAMILY\_V1}      & family gate, 18/18 \\
\bottomrule
\end{tabularx}
\end{center}

The Riemann tensor these contractions are built from is pinned separately.  This
tree carries two implementations of it, and nothing had compared them; a
contraction built against the wrong one would be wrong in a way that every
symmetry test on the other one would miss.  They agree componentwise, and both
satisfy all four algebraic symmetries --- including pair exchange, which is a
theorem rather than a manifest property of the definition and which no test in
this tree had previously covered.

\begin{thebibliography}{9}
\bibitem{Noether1918} E.~Noether,
  \emph{Invariante Variationsprobleme}, Nachr.\ Ges.\ Wiss.\ G\"ottingen,
  Math.-Phys.\ Kl.\ (1918) 235.
\bibitem{Bach1921} R.~Bach,
  \emph{Zur Weylschen Relativit\"atstheorie und der Weylschen Erweiterung des
  Kr\"ummungstensorbegriffs}, Math.\ Z.\ \textbf{9} (1921) 110.
\bibitem{HenneauxTeitelboim} M.~Henneaux and C.~Teitelboim,
  \emph{Quantization of Gauge Systems}, Princeton University Press (1992).
\bibitem{BarnichBrandtHenneaux} G.~Barnich, F.~Brandt and M.~Henneaux,
  \emph{Local BRST cohomology in gauge theories}, Phys.\ Rept.\ \textbf{338}
  (2000) 439.
\bibitem{HobsonBarkerLasenby2026} M.~P.~Hobson, W.~E.~V.~Barker and
  A.~N.~Lasenby, \emph{Bessel--Hagen currents for linearised Weyl-squared
  gravity}, arXiv:2608.02912 (2026), submitted to Phys.\ Rev.\ D.
\bibitem{BoulangerHenneaux2001} N.~Boulanger and M.~Henneaux,
  \emph{A derivation of Weyl gravity}, Annalen Phys.\ \textbf{10} (2001) 935,
  \href{https://arxiv.org/abs/hep-th/0106065}{arXiv:hep-th/0106065}.
\end{thebibliography}

\end{document}
